\hypertarget{evaluate-the-claim-behind-the-benchmark-score}{%
\subsection{Evaluate the claim behind the benchmark
score}\label{evaluate-the-claim-behind-the-benchmark-score}}

\label{exec13:census-evaluation}

CausalEval places diverse language tasks in one comparative survey,
including commonsense cause selection, correlation-to-causation
questions, graph-based inference and judgments about described events.
This is useful navigation, but the common label ``causal reasoning''
does not make their outcomes interchangeable. Correctly classifying a
relation mentioned in a sentence, recovering a graph from observations
and identifying a policy value test different information and
assumptions. \cite{exec13_census_causaleval}

\begin{longtable}[]{@{}lll@{}}
\toprule
\begin{minipage}[b]{0.30\columnwidth}\raggedright
Evidence object\strut
\end{minipage} & \begin{minipage}[b]{0.30\columnwidth}\raggedright
What the inspected source measures\strut
\end{minipage} & \begin{minipage}[b]{0.30\columnwidth}\raggedright
Additional evidence needed for the proposed causal role\strut
\end{minipage}\tabularnewline
\midrule
\endhead
\begin{minipage}[t]{0.30\columnwidth}\raggedright
Susanti--Holsmoelle relation extraction\strut
\end{minipage} & \begin{minipage}[t]{0.30\columnwidth}\raggedright
Binary relations between entities in text\strut
\end{minipage} & \begin{minipage}[t]{0.30\columnwidth}\raggedright
Joint graph constraints, target-population meaning and effect
identification\strut
\end{minipage}\tabularnewline
\begin{minipage}[t]{0.30\columnwidth}\raggedright
Sng and colleagues' hallucination study\strut
\end{minipage} & \begin{minipage}[t]{0.30\columnwidth}\raggedright
Consistency of causal ratings, including a supplied mocked graph\strut
\end{minipage} & \begin{minipage}[t]{0.30\columnwidth}\raggedright
Accuracy against independently justified causal truth\strut
\end{minipage}\tabularnewline
\begin{minipage}[t]{0.30\columnwidth}\raggedright
Augur\strut
\end{minipage} & \begin{minipage}[t]{0.30\columnwidth}\raggedright
Trend prediction and explanation consistency\strut
\end{minipage} & \begin{minipage}[t]{0.30\columnwidth}\raggedright
Intervention recovery under declared temporal and confounding
assumptions\strut
\end{minipage}\tabularnewline
\begin{minipage}[t]{0.30\columnwidth}\raggedright
ORCA and CausalPulse\strut
\end{minipage} & \begin{minipage}[t]{0.30\columnwidth}\raggedright
Workflow outcomes plus distinct scientific metrics\strut
\end{minipage} & \begin{minipage}[t]{0.30\columnwidth}\raggedright
Defined oracles, retained failures and matched target-level
evaluation\strut
\end{minipage}\tabularnewline
\begin{minipage}[t]{0.30\columnwidth}\raggedright
CausalSmith\strut
\end{minipage} & \begin{minipage}[t]{0.30\columnwidth}\raggedright
Reported proof checks and internal run dispositions\strut
\end{minipage} & \begin{minipage}[t]{0.30\columnwidth}\raggedright
Statement fidelity, independent scientific significance and target
assumptions\strut
\end{minipage}\tabularnewline
\bottomrule
\end{longtable}

\hypertarget{relation-extraction-is-a-serious-baseline-with-a-narrower-claim}{%
\subsubsection{Relation extraction is a serious baseline, with a
narrower
claim}\label{relation-extraction-is-a-serious-baseline-with-a-narrower-claim}}

Susanti and Holsmoelle compare prompting with fine-tuned models on four
relation datasets. On DDI, Table 1 reports F1 of 67.4 for Single-Prompt
C and 87.9 for PubMedBERT: a 20.5-point difference. The paper reports
five-fold averages and explicitly limits its evaluation to binary
pairwise causality assessment. It neither demonstrates joint graph
verification nor establishes a universal advantage for fine-tuning under
equal annotation costs. The comparison does establish that a supervised
smaller model deserves consideration when the proposed task is
extraction from a stable document domain.
\cite{exec13_census_graph_validation}

For banking notes, evaluate the extraction role on held-out documents
with independently specified labels and timestamps. Track abstention and
the treatment of ambiguous or contradictory passages. Then separately
evaluate whether extracted statements are admissible graph evidence. A
sentence that says ``the waiver caused repayment'' is evidence of what
the author claimed; it is not, by itself, evidence that assignment was
unconfounded. A useful agent should retain the quotation, its source and
the alternative causal interpretations.

Sng and colleagues illustrate why agreement is a separate target. Their
retrieval experiment supplies a mocked true causal graph, and their
conclusion explicitly says that reducing inconsistent ratings does not
imply accurate causal prediction. A comparison in which one arm receives
such a graph evaluates its use of supplied information. It does not
measure whether retrieval can discover the graph in a bank setting.
Registering both correct and misleading retrieved evidence is necessary
to expose that distinction. \cite{exec13_hallucination}

\hypertarget{task-specific-gains-should-remain-task-specific}{%
\subsubsection{Task-specific gains should remain
task-specific}\label{task-specific-gains-should-remain-task-specific}}

In CausalEval's Table 3, GPT-4o scores 69 on CLADDER with ordinary
chain-of-thought and 76 with FunctionCall, but its e-CARE scores are 90
and 73, respectively. These are author-reported pass@1 results for the
named versions and samples. The gains and losses should remain visible
rather than being replaced by the claim that tools always improve causal
reasoning. The table is also preferable to silently repeating
inconsistent numerical values in the paper's surrounding discussion.
\cite{exec13_census_causaleval}

The study repeats each question five times with greedy decoding.
Repeated calls can reveal response variation, but five calls to the same
question are not five independent causal problems. For illustration,
suppose there are \(N\) independent mechanisms, \(m\) calls per
mechanism, equal marginal score variance \(\sigma^2\), and common
within-mechanism correlation \(\rho\). Then

\[
 \operatorname{Var}(\bar S)
 =\frac{\sigma^2}{Nm}\{1+(m-1)\rho\}.
\]

At \(N=100,m=5,\rho=1\), the effective independent count is 100, not
500. This is a project-derived sampling illustration, not an estimate of
the paper's unreported correlation. The appropriate prototype analysis
clusters calls and paraphrases by underlying mechanism, reports the
number of distinct mechanisms, and uses paired comparisons when methods
see the same cases.

CausalEval also evaluates some expensive models on 100 cases rather than
the usual 500. Its human row combines scores from original benchmark
papers with a study of three graduate students on 100 questions for
remaining datasets. A combined human/model gap therefore is not one
uniform matched experiment. Any investment argument should state the
comparison actually performed, the cost and failure accounting, and the
uncertainty that remains.

\hypertarget{an-affordable-falsifiable-evaluation-sequence}{%
\subsubsection{An affordable, falsifiable evaluation
sequence}\label{an-affordable-falsifiable-evaluation-sequence}}

The proposed sequence starts with fixed causal mechanisms, explicit
public document provenance and deterministic numerical answers. It then
tests formulation, variable timing, graph uncertainty, identification,
execution and reporting separately, followed by their composition.
Baselines receive the same data, permitted documentation and tools.
Model-only reasoning, retrieval, a typed numerical workflow, and a
supervised extractor answer different engineering questions and should
be included where relevant.

The primary denominator includes every attempted case, with scientific
refusals distinguished from crashes and wrong answers. Numerical
accuracy conditional on successful execution is a secondary outcome.
Cost includes failed calls, retrieval, annotation, retries and human
review. Deliberately adverse cases include an accurate quotation
supporting a false causal claim, a timestamp that makes a proposed
control post-treatment, alternative graphs with different effects, a
misleading retrieved graph, and an unsupported edit policy.

Public and generated data can establish that these failure cases are
handled as specified and can compare methods under a known law. They do
not establish the frequency of those failures in a target bank or the
value of changing its real decisions. The next expensive experiment
should address a remaining uncertainty that could change the build or
adoption decision. Until such evidence exists, the proposal can
recommend a bounded prototype and a testable research direction while
preserving uncertainty about live-model performance, source completeness
and bank transport.

\label{exec13:census-evaluation-end}
